\documentclass[11pt]{article}
\usepackage[margin=1in]{geometry}
\usepackage[T1]{fontenc}
\usepackage{xcolor}
\usepackage[authoryear,round]{natbib}
\usepackage[hidelinks]{hyperref}
\setlength{\parindent}{0pt}
\setlength{\parskip}{0.6em}
\setlength{\emergencystretch}{2em}
\begin{document}
\begin{center}
{\Large\bfseries Response to Associate Editor Comment 1}\\[0.4em]
{Knowledge, Retrieval-Augmented Generation, and Modeling Skills}
\end{center}
\textbf{Associate Editor's comment.}
\begin{quote}\itshape
The first one is about the positioning of the paradigm proposed in the paper.
It seems to me that the main message of this paper is that it is important to
have a ``reference'' for each class of problem rather than using a straightforward
agent or fine-tuning approach. Using today's terminology, it would be useful to
have a category-specific ``skill'' as well as ``knowledge'' to deal with domain
specific questions. The authors may want to comment on this and perhaps connect
with some advances in this area.
\end{quote}
% Revised October 6, 2026. Checked against the 0927 notebook.
% Revision locations follow the author's stated manuscript additions.
{\color{blue}
\noindent\textbf{Response:}
We thank the Associate Editor for suggesting this perspective.
We first clarify knowledge, agent skills, and retrieval-augmented
generation (RAG), and then explain how these concepts relate to
LEAN-LLM-OPT's components and methodological positioning.

\medskip
\noindent\textbf{1. Knowledge, skills, and RAG: definitions and relationships.}

\emph{Knowledge.} In this discussion, knowledge refers to the
information and know-how available for completing a task.
Prior work distinguishes declarative knowledge, such as facts
needed to solve a problem, from procedural knowledge, such as
strategies for solving tasks \citep{ae1_li2024metacognitive}.
Thus, knowledge includes both problem-relevant content and
expertise about how to use it; it is not limited to factual
information.

\emph{Skills.} Recent research characterizes agent skills as
reusable procedural artifacts containing instructions, supporting
resources, and applicability conditions
\citep{ae1_zhou2026skillsurvey}. OpenAI describes Agent Skills as
reusable task instructions and supporting files, organized around
a \texttt{SKILL.md} file with optional references, scripts, and
templates \citep{ae1_openai2026skills}. A skill therefore makes a
task method available in a reusable form and can include knowledge
needed to carry out that method. Neither this definition nor the
OpenAI package format requires that skills be learned autonomously
or produce identical actions and outputs for every instance.

\emph{RAG.} RAG retrieves external content and uses it to inform
generation \citep{ae1_lewis2020rag}. OpenAI describes it as
retrieving content to augment an LLM's prompt before generating
an answer \citep{ae1_openai2026accuracy}. It is a mechanism for
bringing relevant external material into generation, rather than
a particular type of knowledge or a format for representing a
task procedure.

\textbf{According to these definitions in OpenAI's official
 documentation, RAG and Agent Skills are distinct concepts.}
Knowledge is the content and know-how available to the agent;
a skill organizes reusable procedural guidance and supporting
resources; RAG supplies retrieved content for generation.
A skill may guide the use of RAG by specifying what information
to retrieve and how to apply it to the task
\citep{ae1_openai2026skilltools}. In this combination, the skill
provides reusable task instructions, while RAG provides the
retrieval-and-generation mechanism.

\medskip
\noindent\textbf{2. Reusable classification and modeling procedures in LEAN-LLM-OPT.}

Viewed through these definitions, LEAN-LLM-OPT organizes agent
behavior around two kinds of \emph{skill-like procedures}:
problem classification, and problem-category-specific modeling
and code generation. Each combines reusable instructions,
reference knowledge, and tool use. The agents execute these
procedures for the current problem.

\emph{Problem classification.}
The classification agent follows a reusable method: identify the
objective, decision variables, and principal constraints; retrieve
relevant labeled reference problems through \texttt{FileQA}; and
use that evidence to determine the problem's structural category.
The guidance specifies how to interpret retrieved labels and
prioritize the current problem's mathematical structure over
surface narrative similarities. This can be understood as a
skill-like classification procedure that uses RAG to obtain
supporting examples. The procedure is reused across problems,
while the retrieved cases and the classification result vary.

\emph{Category-specific modeling and code generation.}
For each predefined problem category, LEAN-LLM-OPT provides
reusable guidance for retrieving reference cases, constructing
formulations, and generating executable code. Reference problems,
expert formulations, and associated data are assembled into
modeling demonstrations. The instructions specify when to call
\texttt{CSVQA}, how to interpret its results, how to define model
components, and what the formulation should contain. The
code-generation stage combines the current formulation with
reference formulations, code examples, and programming instructions
to produce an executable model for the target instance. Together,
these instructions, references, and tool-use routines constitute
a \emph{skill-like procedure} for category-specific modeling and
implementation. For Others and Mixture, shared guidance supports
an abstract modeling plan based on the current query and data
schema, followed by code generation.

These procedures also specify how tool results are checked and
how data-validation failures are handled. In the NRM route, a
model-proposed extraction plan is validated and executed in Python.
If validation fails, the tool returns the complete source data,
and the modeling instructions require the agent to select the
relevant data according to the original query. In the other
standard routes, extracted records are checked against the source
to ensure that identifiers and values are preserved. If this
check fails, the tool returns the original source records.

\emph{Knowledge and RAG within these procedures.}
Ref-Data contains problem descriptions, associated data, expert
formulations, and annotations identifying relevant data; the
implementation also uses reference code. These resources provide
modeling knowledge. They convey both model content and, through
worked demonstrations and code, information about how a model is
constructed. RAG makes relevant cases available to the agents;
the reusable procedural guidance specifies how those cases
should be applied. This use of retrieved examples also connects
to retrieval for in-context learning \citep{ae1_rubin2022epr}.
Reference data support the examples, while target datasets supply
the parameters for the current formulation.

Thus, LEAN-LLM-OPT is functionally close to an agent system that
uses reusable classification and modeling skills together with
RAG. Its current implementation expresses these procedures in
prompts, workflow-construction code, and tools, rather than
separate \texttt{SKILL.md} packages. The reusable method is fixed;
the selected references, detailed modeling plan, tool interactions,
and generated formulation depend on the instance.

\medskip
\noindent\textbf{3. Connection to recent skill research and scope of the framework.}

Agent Workflow Memory induces reusable workflows from task
experience and supplies them to agents in offline or online
settings \citep{ae1_wang2025awm}. In optimization modeling,
OptSkills groups problems by structural archetypes and distills
successful modeling and solving trajectories into reusable
skills, with mechanisms for refinement and expansion
\citep{ae1_yang2026optskills}. These works connect to LEAN-LLM-OPT
through their emphasis on procedural reuse, while also studying
how such procedures can be obtained from experience.

LEAN-LLM-OPT obtains its references and procedural guidance
through design-time curation. They remain fixed during evaluation,
while selected cases, plans, and execution vary with the instance.
The current system neither represents each procedure as an
independently discoverable \texttt{SKILL.md} package nor converts
completed modeling trajectories into new or revised skills in a
persistent library. Skill packaging and skill learning are separate
design choices; the absence of either should not be confused with
the absence of reusable procedural guidance. They are possible
extensions of the current implementation.

\medskip
\noindent\textbf{4. Methodological positioning and relation to fine-tuning.}

We therefore position LEAN-LLM-OPT as a \emph{reference-guided
agentic workflow construction framework} with reusable,
\emph{skill-like classification and modeling procedures}.
RAG supplies reference knowledge for these procedures, and the
agents apply that knowledge through guided reasoning and tool
use to formulate and implement the current optimization model.

These resources could also be used with a fine-tuned model,
consistent with OpenAI's discussion of combining RAG and
fine-tuning \citep{ae1_openai2026accuracy}. LEAN-LLM-OPT makes
curated expertise available at inference time without requiring
an additional task-specific parameter update. Our positioning
therefore does not imply that references replace agents or are
universally preferable to fine-tuning.

\medskip
\noindent\textbf{5. Corresponding manuscript revisions.}

We have added \emph{Knowledge and Skills in LLM Agents} to
Section~1.2 (\emph{Related Works}) to explain the definitions
and relationships of knowledge, skills, and RAG, together with
relevant literature on procedural reuse and skill learning.
We have also added \emph{Connecting Reference Knowledge and
Modeling Skills} to Section~3.3 (\emph{Novelties in Our Design
and Key Messages}) to explain how LEAN-LLM-OPT combines reference
cases, category-conditioned procedural guidance, and target data,
and to clarify the distinction between its skill-like modeling
procedures, standalone skill packaging, and trajectory-based
skill learning.
}

\begin{thebibliography}{99}
\bibitem[Li et~al.(2024)]{ae1_li2024metacognitive}
Li, Z., et~al. (2024). \emph{Meta-Cognitive Analysis: Evaluating Declarative and Procedural Knowledge in Datasets and Large Language Models.} LREC-COLING, 11222--11228. \href{https://aclanthology.org/2024.lrec-main.980/}{Source}.

\bibitem[Lewis et~al.(2020)]{ae1_lewis2020rag}
Lewis, P., et~al. (2020). \emph{Retrieval-Augmented Generation for Knowledge-Intensive NLP Tasks.} NeurIPS 33, 9459--9474. \href{https://proceedings.neurips.cc/paper/2020/hash/6b493230205f780e1bc26945df7481e5-Abstract.html}{Source}.

\bibitem[Rubin et~al.(2022)]{ae1_rubin2022epr}
Rubin, O., Herzig, J., and Berant, J. (2022). \emph{Learning To Retrieve Prompts for In-Context Learning.} NAACL-HLT, 2655--2671. \href{https://aclanthology.org/2022.naacl-main.191/}{Source}.

\bibitem[Zhou et~al.(2026)]{ae1_zhou2026skillsurvey}
Zhou, Y., et~al. (2026). \emph{A Comprehensive Survey on Agent Skills: Taxonomy, Techniques, and Applications.} arXiv:2605.07358v3. \href{https://arxiv.org/abs/2605.07358v3}{Source}.

\bibitem[OpenAI(2026a)]{ae1_openai2026accuracy}
OpenAI (2026a). \emph{Optimizing LLM Accuracy.} API documentation, especially the sections on retrieval-augmented generation and combining RAG with fine-tuning. Accessed October 6, 2026. \href{https://developers.openai.com/api/docs/guides/optimizing-llm-accuracy}{Source}.

\bibitem[OpenAI(2026b)]{ae1_openai2026skills}
OpenAI (2026b). \emph{Skills.} API documentation. Accessed October 6, 2026. \href{https://developers.openai.com/api/docs/guides/tools-skills}{Source}.

\bibitem[OpenAI(2026c)]{ae1_openai2026skilltools}
OpenAI (2026c). \emph{Skills: Understand How Skills Add Repeatable Workflows Around MCP Tools.} Plugins documentation. Accessed October 6, 2026. \href{https://developers.openai.com/plugins/concepts/skills}{Source}.

\bibitem[Wang et~al.(2025)]{ae1_wang2025awm}
Wang, Z. Z., Mao, J., Fried, D., and Neubig, G. (2025). \emph{Agent Workflow Memory.} ICML, PMLR 267:63897--63911. \href{https://proceedings.mlr.press/v267/wang25bx.html}{Source}.

\bibitem[Yang et~al.(2026)]{ae1_yang2026optskills}
Yang, H., et~al. (2026). \emph{OptSkills: Learning Generalizable Optimization Skills from Problem Archetypes via Cluster-Based Distillation.} arXiv:2605.29829v2. Authors report acceptance at Findings of EMNLP 2026. \href{https://arxiv.org/abs/2605.29829v2}{Source}.

\end{thebibliography}
\end{document}
